\begin{figure}
  \begin{minipage}[t]{0.49\textwidth}
\begin{lstlisting}[language=CSharp]
using System;
using System.Drawing;
using System.IO;
using System.Threading.Tasks;
using System.Windows.Forms;

static class Program {
  [STAThread]
  static void Main() {
      Application.Run(new TextEdit());
  }
}

public class TextEdit : Form {
  public TextEdit() { 
    /* */ 
    InitializeServices()
  }

  private void InitializeServices() {
    /* elided */
    _autoSaveController.StartAutoSaveTimer();
  }
}
\end{lstlisting}
  \end{minipage}\hfill%
  \begin{minipage}[t]{0.49\textwidth}
\begin{lstlisting}[
  language=CSharp,%
  firstnumber=25,%
  linebackgroundcolor={\lineEq{42},\lineEq{43}},%
]
public class AutoSaveController {
  private DocumentManager _doc;
  private DiskIO _disk;
  private Action<string> _uiNotify;
  private Forms.Timer _timer;

  public AutoSaveController(
    DocumentManager doc, DiskIO disk,
    Action<string> uiNotify) {
      /* setup elided */
      _timer.Tick += AutoSaveTimer_Tick;
  }

  private async void AutoSaveTimer_Tick(
    object sender, EventArgs e
  ) {
    _uiNotify("Auto-saving...");
    await _disk.SaveFileAsync(
      )_doc.FilePath, _doc.Content);
    _doc.MarkSaved();
    _uiNotify(
      $"Saved {DateTime.Now:T}");
  }
}
\end{lstlisting}
  \end{minipage}
  \cprotect\caption{The above \CSharp{} program is a simple TextEdit application built with Windows Forms. The editor runs a background task that saves the open files periodically. The method \code{AutoSaveTimer_Tick} is \textit{async void}, meaning that any exception raised is placed on the executing thread rather than on an enclosing \code{Task} object. If the highlighted disk save operation raises an exception due to, e.g., insufficient disk space, the entire application crashes! Outside of the method \code{AutoSaveTime_Tick}, a try catch block cannot be used to catch the raised exception.}
  \label{fig:wat:csharp-exn}
\end{figure}
